\section{Related Work}\label{sec:related}

\paragraph{Spectral control in implicit representations.}
SIREN \citep{sitzmann2020siren} sets the frequency scale $\omega_0$ by hand;
Fourier features \citep{tancik2020fourier} expose the bandwidth $\sigma$ and
connect it to the neural tangent kernel \citep{jacot2018ntk}; FINER
\citep{liu2024finer} widens the reachable frequency set. A recent group
selects these parameters from the spectrum of the target signal: FreSh
\citep{kania2025fresh} matches the spectrum of the untrained model to the
target's; WINNER \citep{winner2025} sets Gaussian weight-perturbation scales
from the target spectral centroid; SEC \citep{sec2026} retrieves
grid-search-derived frequency parameters using centroid-based calibration.
WIEN-INR \citep{wieninr} uses wavelet-band subnetworks and a finest-detail
enhancement module. These studies include signal fitting, and FreSh and
WINNER also evaluate inverse tasks. We focus on observation- and
noise-conditioned linear references, mismatch costs, and the empirical
reliability of guidance under degraded observations. Spectral bias
\citep{rahaman2019spectral} explains what an untuned network prefers; our covariance-conditioned reference describes recovery under a specified signal law and observation model.

\paragraph{Arbitrary-scale super-resolution.}
LIIF \citep{chen2021liif} decodes a local latent at continuous coordinates
with a cell (footprint) input; LTE \citep{lee2022lte} adds a local dominant
frequency estimator; MetaSR \citep{hu2019metasr}, SRNO \citep{wei2023srno},
CiaoSR \citep{cao2023ciaosr} and HIIF \citep{jiang2025hiif} refine the
decoder; the SC-INR research implementation examined here replaces the cell input with an
analytic footprint (sinc) readout and includes scale-consistency components; their benefits are evaluated separately from the footprint. Gaussian-splatting decoders
\citep{chen2025gsasr} and heat-field decoders with analytic anti-aliasing
\citep{becker2025thera} form a newer family. A controlled re-training study
\citep{nasir2026empirical} examines the influence of decoder architectures and training recipes. Low-pass behaviour of SR networks has been measured with a global
scalar \citep{fsds2024}; we measure recovery per band defined by the sampling
theorem and set it against a covariance-conditioned linear reference. We examine these empirical similarities and their limitations.

\paragraph{Multi-frame and sparse reconstruction.}
SuperF \citep{superf2026} fits a Fourier-feature network to multiple shifted
low-resolution frames and tunes $\sigma$; Deep Image Prior
\citep{ulyanov2018dip} reconstructs without training data; kriging/GP
regression \citep{rasmussen2006gp} inserts the spectrum directly as a prior.
These tasks serve as linear and inverse-problem checks of the diagnostic, separate from the ASISR readout analysis.

\paragraph{Classical theory and per-band diagnostics.}
LMMSE/Wiener estimation and aliasing analysis of sampling are textbook
\citep{kay1993estimation}. Per-frequency agreement measures such as
Fourier ring correlation in microscopy \citep{vanheel2005frc} and
scale-decomposed skill scores in forecast verification are close relatives of
our per-band recovery score. We use the score to compare scale-dependent ASISR behaviour, rather than claiming a new information measure or a general spectral-selection theorem.
